% Generated draft: requires mathematical and source review.
\documentclass[11pt]{article}
\usepackage[margin=25mm]{geometry}
\usepackage{fontspec}
\setmainfont{Latin Modern Roman}
\usepackage{amsmath,amssymb,mathtools,unicode-math}
\setmathfont{Latin Modern Math}
\usepackage{enumitem,needspace,xurl,hyperref}
\setlength{\parindent}{0pt}
\setlength{\parskip}{5pt}
\setlength{\emergencystretch}{3em}
\newcommand{\N}{\mathbb{N}}
\newcommand{\Z}{\mathbb{Z}}
\newcommand{\Q}{\mathbb{Q}}
\newcommand{\R}{\mathbb{R}}
\newcommand{\C}{\mathbb{C}}
\newcommand{\F}{\mathbb{F}}
\newcommand{\A}{\mathbb{A}}
\newcommand{\PP}{\mathbb{P}}
\newcommand{\E}{\mathbb{E}}
\newcommand{\eps}{\varepsilon}
\newcommand{\dd}{\,\mathrm{d}}
\begin{document}
\section*{UnsolvedMath ID: 30000567}
\textbf{Source problem code:} catalogue entry
\subsection{Definitions and mathematical statement}
\textbf{Disjoint-hypercyclicity question.} Does every separable infinite-dimensional Banach space support a pair, or a larger family, of disjoint-hypercyclic operators? For a Banach space \(X\) over the intended scalar field and bounded linear operators \(T_1,\ldots,T_m\in\mathcal B(X)\), \(m\ge2\), disjoint hypercyclicity means that some \(x\in X\) has dense diagonal orbit\[\{(T_1^nx,\ldots,T_m^nx):n\in\Z_{\ge0}\}\quad\text{in }X^m\] with the product norm topology. Equivalently, every product of nonempty open subsets of \(X\) is hit by this orbit at a common iterate. The same initial vector is used in every coordinate.\par The pair question is \(\forall X\,\exists T_1,T_2\,\exists x\) with that density property. A finite-family version asks it for each prescribed \(m\ge2\); this is not automatically a request for one countably infinite family. For an infinite index set, a topology on the product and the intended family size must be specified. The source's \emph{larger family} does not fix them. No commutativity, invertibility, or common operator criterion is imposed.\par This condition is stronger than merely finding a vector hypercyclic for each operator separately, since all target coordinates must be approximated at the same times. It is also more specific than hypercyclicity of the direct sum with an arbitrary initial tuple: the initial tuple here is \((x,\ldots,x)\). The source does not specify real versus complex scalars; these versions should be stated separately or with the intended field fixed.\par The supplied triage checked 2026-08-21 labels the item open and reports no located authoritative resolution, with the workshop report as its only reference. No precise later theorem or hypothesis comparison is supplied. This dated status assertion is retained without independently claiming that the existence problem remains unresolved today.
\subsection{Short English statement}
Can every separable infinite-dimensional Banach space carry two operators whose synchronized orbits from one common vector are dense in the product space? Larger finite and infinite families require separate quantifiers. Individual or common hypercyclicity alone is weaker.
\subsection{Sources}
\begin{itemize}\item Local UnsolvedMath registry, permanent ID 30000567, code catalogue entry; \emph{Disjoint-Hypercyclic Operators on Banach Spaces}, Oberwolfach Reports Open Problems.\item \emph{Mini-Workshop: Hypercyclicity and Linear Chaos (2007)}, Oberwolfach Reports, DOI 10.4171/owr/2006/37, pp. 2227--2276. \url{https://doi.org/10.4171/owr/2006/37}. Proposed year and year tag: 2006; no proposer named.\item The same full report is the sole supplied literature source. No separate authoritative operator-construction theorem is supplied.\item Literature triage checked 2026-08-21, status \texttt{open}; record updated 2026-09-16. Statement status \texttt{exact}, reported retained byte-for-byte. Verify the dated unresolved-status claim independently.\end{itemize}
\end{document}
